\documentclass[11pt]{article}
\usepackage[margin=1in]{geometry}
\usepackage{booktabs}
\usepackage{hyperref}
\usepackage{url}
\usepackage{tabularx}
\usepackage{array}
\usepackage{longtable}
\usepackage{enumitem}
\hypersetup{colorlinks=true,linkcolor=blue,urlcolor=blue}

\title{Meeting Draft: A Stepwise Framework for\\
Source-Conditioned Synthetic Discharge-Summary Generation}
\author{}
\date{}

\begin{document}
\maketitle

%=============================================================================
\section*{Executive summary}
%=============================================================================

\textbf{Clinical problem.}
Restricted access to large-scale discharge summaries limits reproducible clinical NLP, fair benchmarking, and transparent model development across institutions and patient subgroups.

\textbf{Proposed contribution.}
We are building a \textbf{stepwise, leakage-aware pipeline} for \textbf{source-conditioned synthetic discharge-summary generation}. The project covers baseline generation, coverage-guided enrichment, optional LLM-based editing, and downstream clinical NLP evaluation.

\textbf{Main idea.}
\begin{quote}
An embedding-conditioned language model (ELM) can decode clinically meaningful discharge summaries from held-out real note embeddings. On top of that baseline, targeted embedding-space steering may expand synthetic-note coverage into underrepresented real-note regions, and later editing/evaluation stages may improve note quality and clinical usefulness.
\end{quote}

\textbf{Project structure.}
\begin{itemize}[nosep]
  \item \textbf{Phase 1:} build the auditable vanilla baseline and map the real embedding manifold,
  \item \textbf{Phase 2:} enrich under-covered regions with controlled steering and matched baselines,
  \item \textbf{Phase 3:} improve note quality with LLM-based editing if needed, and test downstream clinical NLP utility.
\end{itemize}

\textbf{How to read the plan.}
Each stage adds one concrete capability to the project. The plan is meant to be directly executable: follow the stages in order, check the gate at each stage, then move to the next block.

%=============================================================================
\section{Step-by-step project roadmap}
%=============================================================================

\begin{longtable}{p{2.0cm} p{3.8cm} p{3.4cm} p{4.1cm} p{1.8cm}}
\toprule
\textbf{Stage} & \textbf{What we do} & \textbf{Functionality added} & \textbf{Why it matters} & \textbf{Current status} \\
\midrule
1. Cohort + leakage &
Prepare the filtered MIMIC-IV discharge-summary cohort, align HADM/note ids, and regenerate the filtered-aligned split manifest. &
Creates the official dataset and row-level provenance needed by every later run. &
Without exact row alignment and leakage flags, no generation result is scientifically auditable. &
\textbf{Done} \\

2. Vanilla ELM baseline &
Decode held-out filtered test embeddings with ELM and write a structured generation manifest. &
Establishes the first source-conditioned synthetic-note baseline with exact row identity. &
This is the anchor condition for all later comparisons. If vanilla itself is not usable, steering is irrelevant. &
\textbf{Done} \\

3. Vanilla audit &
Re-embed generated notes, run manifest integrity checks, quality checks, faithfulness metrics, and privacy/memorization screens. &
Turns raw generated notes into an evaluated baseline with PASS/CAUTION/FAIL status. &
This is the technical gate before we use synthetic notes for any coverage or enrichment claim. &
\textbf{Done} \\

4. Real-manifold mapping &
Map the full filtered real embedding space across train/dev/test; identify dense, sparse, and subgroup-skewed regions. &
Builds the empirical coverage target in the original 1024-d BGE space. &
This tells us where representation gaps are and prevents arbitrary steering targets. &
\textbf{Done} \\

5. Held-out coverage validation &
Compare held-out real test embeddings with vanilla synthetic embeddings, stratified by patient-disjoint vs.\ overlap rows. &
Measures how well vanilla synthetic notes cover the real held-out manifold. &
This shows what vanilla already covers well and what remains under-covered before any intervention. &
\textbf{Done} \\

6. Steering redesign (Phase 2b) &
Build more target-specific local steering directions, generate steered notes, and compare them against matched vanilla and norm-matched random shift. &
Tests whether embedding-space steering can enrich specific under-covered regions after decoding. &
This is the key methodological question: can we move beyond generic perturbation and create targeted enrichment? &
\textbf{Active} \\

7. Utility extension &
Only after a convincing steering win, evaluate whether privacy-audited synthetic notes help downstream clinical NLP such as NER. &
Links generation quality and coverage to a real clinical NLP endpoint. &
This is how the project moves from methodology to practical clinical value. &
\textbf{Later} \\
\bottomrule
\end{longtable}

\noindent\textbf{Stage gate logic.}
\begin{itemize}[nosep]
  \item We do \textbf{not} start steering claims until vanilla generation passes audit.
  \item We do \textbf{not} scale a steering method unless it beats both matched vanilla and norm-matched random shift.
  \item We do \textbf{not} start LLM editor or downstream NER until we have a convincing synthetic-note generation condition worth carrying forward.
\end{itemize}

%=============================================================================
\section{Implementation snapshot}
%=============================================================================
\subsection{Pipeline overview}

\begin{center}
\small
\textbf{Real note}
$\rightarrow$ \textbf{1024-d embedding}
$\rightarrow$ \textbf{ELM generation}
$\rightarrow$ \textbf{synthetic note}
$\rightarrow$ \textbf{audit}
$\rightarrow$ \textbf{coverage-guided steering}
$\rightarrow$ \textbf{optional LLM editor}
$\rightarrow$ \textbf{downstream evaluation}
\end{center}

\subsection{Current cohort status}

\begin{itemize}[nosep]
  \item MIMIC-IV discharge summaries, note/HADM aligned: $\sim$332k notes embedded
  \item Full split before long-sequence filtering: train 265,434; dev 33,179; test 33,180
  \item Filtered ELM cohort used by the model: train 262,895; dev 32,847; test 32,843
  \item Rows removed by long-sequence filtering: 2,539 train; 332 dev; 337 test (3,208 total)
  \item ELM training complete: \texttt{checkpoint-8215} (Apr 2026)
  \item Pre-filter and filtered-aligned leakage audits complete
  \item Canonical filtered-aligned split manifest now matches \texttt{encoded\_\*\_filtered} and is the official join target for generation and audit
  \item Vanilla held-out generation complete with structured manifest: 32,843 / 32,843 notes generated on the filtered test split
  \item Vanilla audit complete: \textbf{PASS}; median source$\rightarrow$synthetic cosine 0.7892 overall and 0.7864 on the patient-disjoint subset; empty-output rate 0; collapse rate 9.23\%
  \item Whole-filtered-real coverage discovery complete on 328,585 filtered embeddings across train/dev/test
  \item Held-out real-vs.-vanilla coverage baseline complete; vanilla covers all 10/10 low-density held-out test clusters and reaches real$\rightarrow$synthetic coverage 0.628 overall (0.609 on patient-disjoint rows)
  \item Subgroup metadata table is now built for age bin, sex/gender, race/ethnicity, insurance, admission type, service, LOS bin, and ICU flag
\end{itemize}

\subsection{Why the manifest matters}

The project uses two linked provenance layers:
\begin{enumerate}[nosep]
  \item a \textbf{whole-cohort source manifest} for stable embedding/source provenance,
  \item a \textbf{filtered-aligned split manifest} that exactly matches \texttt{encoded\_\*\_filtered}.
\end{enumerate}

The generation manifest is required because it lets us:
\begin{itemize}[nosep]
  \item keep exact row alignment between real embeddings, generated notes, and metadata,
  \item stratify all results by leakage flags,
  \item compare later conditions on the same source anchors,
  \item carry provenance into steering, editing, and downstream evaluation.
\end{itemize}

%=============================================================================
\section{Research questions}
%=============================================================================

\paragraph{RQ1 (Phase 1). Faithfulness and quality.}
Does vanilla, source-conditioned ELM generation produce readable discharge summaries that preserve source-note semantics on held-out embeddings?

\paragraph{RQ2 (Phase 1). Coverage baseline.}
What regions of the \textbf{full filtered real-note manifold} are dense, sparse, or subgroup-skewed, and which of those regions are under-covered by vanilla synthetic generation on held-out evaluation?

\paragraph{RQ3 (Phase 2). Audited enrichment.}
Can embedding-space steering improve coverage of a clinically meaningful under-covered \textbf{local basin} beyond matched vanilla and norm-matched random perturbation, without unacceptable faithfulness loss or increased privacy risk?

\paragraph{RQ4 (Phase 3). Clinical utility.}
Do privacy-audited synthetic notes improve downstream clinical entity extraction relative to real-note-only development?

%=============================================================================
\section{Core evaluation axes}
%=============================================================================

All outputs are evaluated on the same core axes:

\begin{itemize}[nosep]
  \item \textbf{Faithfulness:} source$\leftrightarrow$synthetic cosine, retrieval rank, neighborhood recovery, and sampled content checks,
  \item \textbf{Primary analysis space:} the full normalized 1024-d BGE embedding space,
  \item \textbf{Coverage:} cluster occupancy, low-density coverage, subgroup coverage, and real-vs.-synthetic distances,
  \item \textbf{Quality:} note length, structure, repetition/collapse, and manual review,
  \item \textbf{Privacy:} duplicate detection, overlap checks, nearest-neighbor screens, and PHI-like pattern scans,
  \item \textbf{Leakage:} note/HADM/patient overlap reporting and patient-disjoint sensitivity analysis,
  \item \textbf{Utility (Phase 3):} downstream clinical NLP performance with and without synthetic-note augmentation.
\end{itemize}

PCA/UMAP/t-SNE are used for visualization only; coverage claims come from the original high-dimensional embedding space.

%=============================================================================
\section{Stage details and decision logic}
%=============================================================================

\subsection{Stage 1: Build the auditable baseline}

Goal: prove that source-conditioned ELM generation is technically valid and scientifically reportable.

\textbf{What we do.}
\begin{enumerate}[nosep]
  \item finalize the filtered cohort and filtered-aligned leakage manifest,
  \item run vanilla held-out generation with row-level manifest writing,
  \item audit quality, faithfulness, privacy risk, and patient-disjoint sensitivity.
\end{enumerate}

\textbf{Functionality.}
This stage gives us the official baseline corpus, manifest, and audit outputs.

\textbf{Why it is important.}
If this baseline is weak or non-reproducible, later enrichment would be uninterpretable.

\textbf{Decision gate.}
Proceed only if the vanilla run is manifest-clean, readable, non-collapsed, privacy-risk-screened, and acceptable on patient-disjoint rows.

\subsection{Stage 2: Map the real manifold and define the coverage target}

Goal: understand the structure of the real filtered embedding space before trying to ``improve'' anything.

\textbf{What we do.}
\begin{enumerate}[nosep]
  \item cluster the full filtered real embeddings across train/dev/test,
  \item summarize sparse and dense regions,
  \item join subgroup metadata to identify clinically meaningful skews,
  \item compare held-out real test rows against vanilla synthetic coverage.
\end{enumerate}

\textbf{Functionality.}
This stage builds the coverage map that later steering methods must target.

\textbf{Why it is important.}
It turns the project from ``make more synthetic notes'' into ``fill empirically under-covered clinical regions.''

\textbf{Decision gate.}
Proceed to steering only after we know which regions are under-covered and can justify why they matter clinically.

\subsection{Stage 3: Test whether steering adds targeted value}

Goal: determine whether target-specific steering can enrich specific real-manifold regions after decoding.
\textbf{Current conclusion.}
The first global-axis pilots, local additive pilots, projected-boundary pilots, and lightweight decoder-adaptation pilots were scientifically useful, but none justified inference-time scale-up. The main lesson is now clear:
\begin{quote}
good pre-decode geometry and acceptable note quality do \emph{not} guarantee that decoded notes will occupy the intended sparse real-note region.
\end{quote}
The decoder-feasibility audit confirmed the limiting mechanism: even real target-basin embeddings weakly retained basin identity after ELM decode and BGE re-embedding. Therefore, the active method is no longer another input-space steering family. It is \textbf{closed-loop output-space enrichment}: generate several vanilla candidates from a real target anchor, re-embed each output, and retain only candidates that satisfy target-region, faithfulness, quality, privacy-risk, and diversity gates.

\textbf{Current closed-loop result.}
Using 256 held-out target-basin anchors, eight candidates per anchor, default prompting, and an 8192-token output budget produced 2,048 candidates and 106 strict accepted notes (5.18\%). This is a substantial improvement over the earlier 14-note strict run (0.68\%), largely because completion and structure improved at the longer output budget. The selected notes improve exact pooled-basin landing and automated quality relative to fair one-draw matched vanilla. They occupy the pooled local basin around clusters 9/17/29/45, while exact cluster-29 occupancy remains a secondary diagnostic. However, this result is now interpreted as \textbf{output-space geometry-selection feasibility}, not as clinically usable source-conditioned note synthesis.

\textbf{Validation constraint.}
Automated acceptance is not enough for clinical use. The completed row-level review found missing or non-substantive sections in 205/206 blinded notes, frequent medication-reconciliation and physiologic-value errors, and severe source-anchor loss in the source-paired sample: principal diagnosis was not preserved in 67--68/68 cases, every synthetic variant contained an unsupported major claim and omitted a critical source fact, and no meaningful rescue appeared for exact-pool versus centroid-only selection. Thus, the selected notes can be geometrically target-aligned while representing a different clinical encounter. We do \textbf{not} scale candidate count, claim clinical enrichment, start source-free LLM editing, or begin downstream NER from this condition.

\textbf{Definitive selector analysis.}
The completed manifests now support three complementary analyses without new generation. First, selected notes improve exact pooled-basin landing over one-draw matched vanilla on the same 106 anchors (32.1\% versus 15.1\%; paired risk difference +17.0 percentage points, exact McNemar $p=0.0079$). Second, the all-anchor operational yield is 34 usable exact-basin notes (13.3\%), 72 centroid-only notes (28.1\%), and 150 anchors without a final accepted note (58.6\%). Third, a random-one-of-eight equal-compute baseline has higher \emph{raw} exact-basin landing (about 21.5\%) but very low clinically-rule-passing exact yield (about 1.75\%), whereas closed-loop selection yields 13.3\% usable exact-basin notes. Thus the primary selector endpoint is \textbf{usable exact pooled-basin enrichment}; raw exact landing is retained as a secondary geometry diagnostic because it can favor collapsed or clinically unusable candidates.

\textbf{Latest refinement from the cluster-29 pilot.}
The most recent upper-bound and basin-margin analyses show that decoded notes are usually not flying to an unrelated part of the manifold. Instead, they remain inside a \textbf{nearby local basin} around cluster 29, especially near clusters 9, 17, and 45. This changes the scientific framing:
\begin{itemize}[nosep]
  \item exact single-cluster occupancy should be treated as a \textbf{diagnostic metric},
  \item the more meaningful Phase 2 target is \textbf{region-level enrichment} of a sparse local basin,
  \item cluster 7 is geometrically nearby but may need to be evaluated separately if it is not clinically coherent with the main basin.
\end{itemize}
In other words, Phase 2 should now ask whether we enrich the intended sparse \emph{region}, not whether every decoded note falls back into one exact KMeans bucket.

\textbf{Current hard-anchor conclusion.}
\begin{itemize}[nosep]
  \item On the true hard-anchor set from the full held-out pool, the current \textbf{linear local basin-margin family fails its strongest test}.
  \item The best hard-anchor sweep setting improves target-vs-competitor margin substantially, but still yields \textbf{zero exact target-cluster wins}.
  \item Therefore the main bottleneck is now clearer: the operator can move embeddings \emph{toward} the local target basin, but it does not cross the \textbf{local competitor boundary} strongly enough to beat clusters 9 and 45 after decoding.
\end{itemize}

\textbf{Ranked redesign options.}
\begin{enumerate}[nosep]
  \item \textbf{Competitor-aware local transport / barycentric crossing.} For each source anchor, build a source-conditioned move toward real target exemplars while explicitly pushing away from the current best competitor. This is the best fit to the observed failure mode because the main problem is \emph{local competition}, not arbitrary off-manifold drift.
  \item \textbf{Competitor-aware basin-margin steering.} Keep the generic basin-margin framework, but redesign the operator so it explicitly optimizes target-cluster margin over clusters 9/45/17 instead of only nudging toward a pooled basin.
  \item \textbf{Decode--reembed correction.} Treat steering as a fixed-point problem. Decode once, re-embed the synthetic note, measure the residual to the intended target region, and run one correction step. This directly addresses the gap between good pre-decode geometry and weak post-decode occupancy.
  \item \textbf{Lightweight projector or adapter refinement.} If inference-only local operators still fail, train a small projector in 1024-d space or fine-tune only the ELM adapter so locally shifted embeddings decode more faithfully without retraining the full language model.
  \item \textbf{Another broad global-axis or plain additive sweep.} Lowest priority. The hard-anchor result suggests this is unlikely to resolve the true bottleneck.
\end{enumerate}

\textbf{Why these options fit our model.}
\begin{itemize}[nosep]
  \item The ELM does \textbf{not} consume retrieved text or multiple control tokens; it consumes one 1024-d embedding through an adapter.
  \item Because of that, methods that keep the edited embedding close to \textbf{real observed embeddings} are more plausible than unconstrained additive moves.
  \item The failure of the projected-boundary pilot suggests that \textbf{preserving cosine to the source alone is not enough}; the edited embedding must also stay on a decoder-compatible manifold.
  \item The failure of broad global directions suggests that steering should be \textbf{conditional on the current source anchor}, not static across all inputs.
\end{itemize}

\textbf{Revised success rule.}
Stage 3 succeeds only if at least one steering condition:
\begin{itemize}[nosep]
  \item first shows a \textbf{pre-decode win on hard anchors} against the local competitor set,
  \item improves patient-disjoint real$\rightarrow$synthetic coverage,
  \item improves basin-level occupancy, local sparse-region coverage, or low-density-cluster coverage in the intended region,
  \item preserves acceptable faithfulness/quality/privacy,
  \item and \textbf{beats both matched vanilla and norm-matched random shift}.
\end{itemize}

\textbf{Immediate next step for the current pilot.}
For cluster 29, the next concrete evaluation block is:
\begin{enumerate}[nosep]
  \item freeze input-space steering and decoder-adaptation results as negative/partial evidence about decoder compatibility,
  \item keep `29/9/17/45` as the current candidate local basin, with `7` evaluated separately,
  \item freeze the 106-note closed-loop run and fair matched-vanilla comparison as a \textbf{negative clinical-faithfulness gate} plus a positive geometry-selection finding,
  \item correct privacy reporting so shared boilerplate is separated from material n-gram copying,
  \item freeze the completed four-case source-fact-conditioned smoke as a method-selection experiment, not a precision estimate: \texttt{checkpoint-8215} fact-only was 4/4 rule-verified passes, whereas raw ELM and both raw-draft correction arms were 0/4,
  \item run a prospective 30-case, one-shot \texttt{checkpoint-8215} fact-only replication with the existing raw ELM outputs as matched baseline and a prespecified 15-case untouched-backbone fact-only control,
  \item require improved diagnosis, procedure, complication, medication-change, disposition, and source-faithfulness preservation before re-embedding clinically passing outputs for geometry or coverage analysis.
\end{enumerate}

\textbf{Current rescue-pilot preparation.}
The completed manual labels have been ingested without copying source text into derived summaries. A reproducible 45-case held-out pilot cohort was selected across exact-pooled and centroid-only routes with 10 patient-disjoint cases. Manual ledger review transformed 406 provisional facts into a completed restricted ledger of 483 rows: 354 verified, 82 corrected, and 47 rejected, with 77 reviewed source-supported additions. All 45 cases now have verified/corrected coverage for the six required fields, and independent source-span validation is 100\% for the 436 usable facts. The local initial backbone and \texttt{checkpoint-8215} text-generation paths both passed deterministic smoke tests. Only four cases currently have reviewer-approved compact generation values; the 30-case replication must prospectively freeze unused cases and complete the same concise generation-ledger review before any generation. No source-derived text may leave approved storage.

\textbf{Functionality.}
This stage turns the project from passive auditing into controlled enrichment.

\textbf{Why it is important.}
This is the project's main methodological innovation, and it only counts as a success if it clearly outperforms the two required controls.

\textbf{Recent literature guiding the redesign.}
\begin{itemize}[nosep]
  \item \textbf{Embedding inversion work} shows that text embeddings preserve substantial semantic content, but inversion often benefits from \emph{iterative correction} rather than a single one-shot decode. This supports decode--reembed correction instead of trusting one edited embedding.
  \item \textbf{Representation-engineering surveys} emphasize that one global steering direction is often too coarse for heterogeneous concepts. This matches our current finding that broad or static directions do not transfer reliably to note-level target occupancy.
  \item \textbf{Targeted intervention papers} show that more selective, conditional steering can outperform untargeted intervention. For our project, that translates to source-conditioned operators rather than one fixed direction for every source note.
  \item \textbf{Conditional activation-steering papers} argue that local, context-aware control is often better than static steering strength. For us, that means alpha and target-neighborhood choice should depend on the source anchor instead of being globally fixed.
\end{itemize}

\textbf{Current status.}
The input-space axis, local-transport, and decoder-adaptation families are now \textbf{negative/partial findings}: they clarify decoder compatibility but do not support inference-time enrichment claims. Closed-loop output selection is a positive feasibility signal for final-output basin landing, but the completed clinical and source-paired reviews fail the source-faithfulness requirement. The four-case source-grounded smoke selected \texttt{checkpoint-8215} fact-only generation, while raw-draft correction failed. The next milestone is a prospective 30-case one-shot factuality replication with a matched raw-ELM baseline and nested untouched-backbone control, not another steering family, geometry optimization, larger candidate pool, or downstream-NLP experiment.

\subsection{Stage 4: Editing and downstream utility}

Goal: improve note quality where needed and test whether the final synthetic corpus helps real clinical NLP tasks.

\textbf{What we add.}
\begin{itemize}[nosep]
  \item a constrained LLM editor for local coherence repair,
  \item downstream clinical entity extraction and augmentation studies,
  \item optional reference-guided coverage analysis from an external cohort if available.
\end{itemize}

\textbf{Why it is important.}
These steps connect the project to the mentor's main interests: better synthetic-note quality and practical clinical NLP value.

\textbf{Important constraint.}
Editing and downstream evaluation only start after we have a synthetic-note condition worth carrying forward. An editor must be source-fact-conditioned and constrained; a source-free editor may polish a clinically wrong encounter. Any edited note must still pass source-alignment and privacy checks.

%=============================================================================
\section{Main risks and mitigation}
%=============================================================================

\begin{longtable}{p{4.2cm} p{5.2cm} p{5.0cm}}
\toprule
\textbf{Risk} & \textbf{Why it matters} & \textbf{Mitigation} \\
\midrule
Vanilla ELM not faithful enough &
  Invalidates all enrichment claims &
  Phase 1 faithfulness/quality gate before CAV or editor \\
CAV adds no value beyond sampling &
  Apparent diversity may be decoder noise &
  Multi-seed vanilla baseline \\
Random shifts match CAV coverage &
  Steering may not be interpretable enrichment &
  Norm-matched random perturbation control \\
Patient overlap inflates results &
  Held-out notes may be too easy for memorization-adjacent patients &
  Patient-disjoint primary sensitivity subset defined from the filtered-aligned manifest \\
Manifest drift after filtering &
  Pre-filter split metadata can stop matching the actual HF datasets used by the model &
  Maintain both source manifest and filtered-aligned manifest; use filtered manifest for all generation/audit joins \\
Synthetic notes memorize training text &
  Privacy and scientific invalidity &
  Duplicate, $n$-gram, NN, and PHI-pattern audits \\
Editor rewrites too much &
  Breaks source conditioning and confounds ELM/CAV claims &
  Post-edit cosine/entity checks; keep editor optional and late \\
Coverage gains are off-manifold &
  Synthetic text may not be clinically meaningful &
  Restrict enrichment to anchored real embeddings; audit entity/clinical content \\
\bottomrule
\end{longtable}

%=============================================================================
\section{Current project status}
%=============================================================================

\noindent\textbf{Current working interpretation.}
The baseline system is established: vanilla generation, manifesting, leakage-aware auditing, and real-manifold mapping are complete. Phase 2 has two distinct findings. First, input-space steering and decoder-adaptation pilots did not reliably preserve sparse-region identity through decode--re-embed, so they are not reliable \emph{input-control} mechanisms. Second, closed-loop output-space selection did produce notes that land in the intended pooled sparse basin better than one-draw matched vanilla, establishing \emph{output-space geometry feasibility}. Its blinded review failed the separate clinical source-faithfulness gate. The active question is therefore whether a \textbf{source-grounded local decoder workflow} can preserve clinical facts while retaining this output-space enrichment capability.

The prospective 30-case fact-only replication resolves the factuality half of this question: all primary checkpoint outputs passed the blinded factual gate. However, the subsequent paired geometry audit shows that factual fact-only notes retain the pooled basin in only 30.0\% of cases overall and 14.3\% of patient-disjoint cases, versus 50.0\% and 100\%, respectively, for the frozen \emph{closed-loop-selected} raw ELM comparator. The latter is not a one-draw vanilla baseline and remains clinically unacceptable. Thus clinical factuality and output-space geometry remain separate requirements. The next bounded bridge is multi-candidate, fact-ledger-conditioned generation followed by final-output basin selection and a new blinded factual review of the selected notes.

The first eight-anchor bridge pilot generated four sampled fact-only candidates per anchor. All 32 candidates were complete and non-empty; 50.0\% landed in the pooled basin, and selecting one output per anchor by final BGE basin score yielded 5/8 (62.5\%) pooled-basin landing, including three patient-disjoint cases. All eight selected notes subsequently passed blinded ledger-grounded factual review (zero unsupported-major claims and zero critical omissions; mean factual faithfulness 4.75/5). This establishes dual-gate feasibility, not scale-up readiness. The next confirmatory step is the existing 30-anchor reviewed cohort with four candidates per anchor, final-output selection, and blinded review of all 30 selected notes.

The 30-anchor confirmation generated 120 sampled fact-only candidates. Candidate pooled-basin landing was 35.0\%; selecting one candidate per anchor by final BGE score yielded 16/30 (53.3\%) pooled-basin landing, split evenly across exact-pooled and centroid-only strata. The selected cohort is technically complete (all outputs non-empty and EOS-terminated) and retains all seven patient-disjoint anchors, although only two patient-disjoint outputs land in the basin. Blinded factual review of the 30 selected notes is now the decisive gate before a second-region pilot or any cohort-scale claim.

The expanded blinded review passed in 28/30 cases (93.3\%), with mean factual faithfulness 4.63/5. Fifteen notes (50.0\% of anchors) jointly passed the factual rule and landed in the pooled basin; all seven patient-disjoint anchors passed factual review, with two dual-gate successes. The two failures were medication-reconciliation errors rather than a broad factuality failure. The frozen result is not retrospectively edited. Future prompts include an explicit ledger-only, non-contradictory medication-list rule. The next stage is a preregistered second-region replication under the same generation, final-output selection, blinded factual review, and leakage-aware reporting protocol.

The first prospective second-region eligibility screen (cluster 20) showed why factual-ledger completeness must precede generation. Of 30 reviewed anchors, only 13 retained all core verified facts after allowing unsupported follow-up to remain absent; the other 17 lacked a verified diagnosis, hospital course, instructions, or disposition. These anchors are excluded rather than filled with invented content. Cluster 20 is therefore frozen as a source-grounding eligibility limitation, and the next preregistered second-region candidate is the clinically distinct surgery-enriched cluster 36.

Cluster 36 exposed a second source-grounding limitation: follow-up sections were placeholder-only and most instruction sections were incomplete fragments. This was confirmed by a restricted aggregate extraction audit rather than assumed from model output. The project will not relabel these fields as present or generate from incomplete ledgers. The next redesign is a source-eligibility pre-screen over the full candidate region, selecting only anchors with complete non-placeholder source support before manual ledger review. Thus source-document completeness becomes an explicit feasibility criterion alongside embedding-region sparsity and leakage status.

A bounded full-note recovery audit was then run before abandoning the region. High-recall heading search recovered substantive follow-up evidence in 13/30 cases and structured admissions provided a disposition candidate in 28/30 cases, but only 3/30 cases had substantive instruction evidence. Therefore evidence recovery is retained as a provenance-aware eligibility tool, not as a license to impute missing patient instructions. Cluster 36 remains unsuitable for a 30-case complete-note replication; diagnoses, procedures, and medication orders remain corroborative rather than automatic discharge facts.

Before selecting another primary target, a reusable source-eligibility pre-screen now evaluates every held-out row in a candidate sparse region. It assigns a complete-note review candidate only when all required fields have substantive source-note evidence, allowing an admissions disposition only as a separately labelled candidate for later verification; it separately labels partial-document candidates and insufficient-evidence rows. The screen exports IDs, evidence classes, routes, and leakage provenance but no source-note text. Manual ledger verification remains mandatory for every complete-note candidate. The next region will be selected from regions with both under-coverage relevance and an adequate complete-note candidate reserve, avoiding repeated manual review of documentation-poor regions.

The first calibrated screen compared clinically distinct candidate regions 11, 16, and 25. Because principal diagnoses and list-form medication/instruction sections can be legitimately short, the screen uses field-specific high-recall criteria and is explicitly a review-triage tool rather than an automatic fact validator. Cluster 16 had 52 complete-note review candidates, including 25 patient-disjoint rows, compared with 19/8 for cluster 11 and 66/6 for cluster 25. A deterministic 30-anchor cluster-16 review cohort (14 patient-disjoint; seed 20260718) is frozen for the next source-ledger verification pass. No generation or geometry claim will be made unless that reviewed cohort clears the factual gate.

For cluster 16, the initial standard section audit encountered early placeholder headings for follow-up despite later substantive evidence elsewhere in the same note. The recovery and ledger tools therefore select the strongest matching heading rather than the first match. The corrected restricted recovery audit identifies substantive follow-up and instruction evidence in all 30 frozen anchors and structured disposition candidates in all 30. This confirms that full-note recovery can improve source-eligibility recall, while manual verification remains the required factual gate before generation.

Manual verification supersedes automatic recovery classification. In the reviewed cluster-16 ledger, only two of 30 cases retained all six fields; in the other 28 cases, the apparent recovery candidate was medication-history or unrelated text rather than a source-supported follow-up plan. The six-field complete-discharge experiment is therefore not pursued for cluster 16. However, every reviewed case retained principal diagnosis, hospital course, discharge medications, disposition, and instructions. We consequently define a separately labelled \\emph{source-supported discharge-transition note} task: these five fields are required, follow-up is included only when explicitly verified, and absent follow-up is omitted rather than generated. This is not a complete-discharge-summary claim. The reviewed 30-case cohort is structurally ready for a bounded transition-note feasibility pilot, followed by the same blinded factual and output-geometry gates used in the first basin.

The first transition-note geometry pilot selected 7/8 outputs into cluster 16 but passed blinded factual review in 6/8 cases; both failures asserted unsupported absence of follow-up. A fresh eight-anchor phrase-heavy prompt guard was therefore tested prospectively and performed worse (3/8 factual passes, four unsupported follow-up claims, and one critical disposition omission), confirming that prompt wording alone is not a reliable factual-control mechanism. The revised dual-gate algorithm applies deterministic factual eligibility before output geometry: candidates mentioning follow-up without a verified follow-up ledger fact are rejected, and only remaining candidates are ranked by final BGE target-region score. The filtered result on the guard pilot is diagnostic only. A new independent pilot is required before any scale-up claim.

The independent cluster-16 confirmation applied that factual preselection prospectively. It rejected 13/32 candidates with unsupported follow-up mentions, retained an eligible candidate for seven of eight anchors, and selected a cluster-16 output for five of the eight frozen anchors. Blinded review of the seven selected notes yielded 6/7 passes, no unsupported-major-claim failures, and one critical omission of a ledger-supported home disposition. All five target-cluster outputs passed, while one of two out-of-basin outputs also passed; this supports separate geometry and clinical gates rather than a causal association claim. Because the prespecified scale-up criterion (at least 7/8 factual passes with no critical omission) was not met, cluster 16 is frozen as a bounded prospective feasibility result and is not scaled. The next replication target is selected from the remaining five-field transition-note eligibility reserve, with cluster 25 providing 317 high-recall candidates including 46 patient-disjoint rows.

The first deterministic 30-anchor cluster-25 draw contained only four patient-disjoint anchors because it used proportional stratification. This preliminary file is not the official replication cohort. The anchor-freezing tool now supports an explicit patient-disjoint minimum; the versioned cluster-25 replication will require at least eight patient-disjoint anchors from the available reserve of 46 before restricted ledger construction. This improves leakage-aware sensitivity reporting without changing the target region, source-evidence contract, or blinded-review gate.

The official cluster-25 transition-note replication cohort is now frozen with 30 anchors, including eight patient-disjoint anchors (seed 20260721). Its restricted provisional ledger contains 287 candidate facts. Hospital-course and instruction candidates are present for all 30 cases. Standard extraction finds placeholder-only follow-up in all cases, whereas high-recall recovery finds substantive follow-up candidates in only 7/30; follow-up consequently remains optional and is omitted unless manual verification supports it. Disposition and other short field candidates remain review targets rather than automatic facts. This establishes source-evidence eligibility for a manual-ledger pass, not permission to generate or a result about clinical quality or geometric enrichment.

Manual cluster-25 ledger review retained 77 verified and 164 corrected compact values while omitting 46 unsafe or unsupported candidates. Twenty-one of the 30 frozen anchors retain all five required transition fields with optional follow-up; nine are blocked rather than imputed. A prompt-safe 21-case ledger and a deterministic eight-case pilot (four patient-disjoint; seed 20260722) are frozen for the same four-candidate fact-only, final-output geometry-selection, and blinded factual-review protocol used for cluster 16. Exact source-string matching is low after reviewer normalization and is retained as a provenance limitation; manual source-evidence review remains the factual authority. This is a bounded second-region feasibility replication, not a cohort-scale enrichment claim.

The cluster-25 pilot provides a sharper dual-gate result. After prospective rejection of candidates that mentioned unsupported follow-up, six of eight frozen anchors retained a selected candidate and all six final outputs landed in cluster 25 (75.0\% frozen-anchor target coverage). Blinded review passed 4/6 selected notes (66.7\%), with no unsupported-major claims but two critical omissions of ledger-supported dispositions (``Home'' and ``Home With Service''). Geometry was therefore compatible with factual generation, but it did not guarantee required transition-field completeness. Candidate-level inspection found a target-landing, follow-up-eligible alternative containing a disposition heading for each failed anchor. The next confirmation will consequently predeclare a \emph{required-field structural eligibility} screen before geometric ranking: for every required ledger field, the candidate must contain a corresponding section heading; unsupported follow-up must still be omitted. This screen is a deterministic completeness guard, not a factuality substitute, and its impact must be assessed prospectively in an independent frozen cohort. Cluster 25 remains a bounded partial feasibility result; it is not ready for scale-up or a clinical-quality claim.

The confirmation cohort must exclude all prior cluster-25 frozen anchors and restore patient-disjoint representation from the larger high-recall eligibility reserve before ledger review. Manual review is used to calibrate and periodically audit the deterministic factual and completeness gates, not as a requirement to review every future synthetic note. Scaling is justified only after an independent cohort demonstrates both adequate final-output target-region yield and a prespecified low rate of unsupported claims and critical required-field omissions in a blinded, leakage-stratified sample. At scale, the same automatic gates will be applied to every candidate, while blinded review remains a stratified quality-control sample and not an output-selection mechanism.

The independent cluster-25 confirmation is now frozen with 12 previously unused anchors (four patient-disjoint; seed 20260724), excluding the original 30-anchor cohort. Manual ledger review retained ten cases with all five required transition fields, including three patient-disjoint cases; two were blocked rather than imputed for missing supported disposition or medication/instruction evidence. The resulting prompt-safe ledger contains 82 compact facts across ten cases. The next prospective test is therefore fixed at four sampled fact-only candidates per cleared case, followed by BGE re-embedding, predeclared follow-up and required-field eligibility screens, final-output geometry selection, and blinded review. It is a 40-candidate confirmation, not scale-up.

The 40-candidate confirmation retained 22 structurally eligible candidates and seven selected outputs from the ten frozen cases; five of ten frozen anchors yielded a final cluster-25 output. Blinded review of all seven selected outputs passed in 6/7 cases, with zero unsupported-major claims and one critical omission of a source-supported antibiotic/home-service transition plan. Thus required-field headings improve structural completeness but do not ensure within-section coverage of every material transition detail. This clears a controlled moderate-scale experiment with the same rejection accounting and blinded stratified quality control, but not full-cohort production. The operational claim is a quality-gated synthetic transition-note subset with transparent yield, not clinical equivalence to the original discharge summaries.

The independent 24-anchor cluster-25 moderate-scale validation generated 96 candidates, retained 20 anchors after predeclared cap, follow-up, and required-heading eligibility screens, and selected a final cluster-25 output for 12/24 frozen anchors. Blinded review of all 20 selected outputs passed in 19/20 cases (95.0\%), with no unsupported-major claims, one medication-reconciliation omission, mean factual faithfulness 4.55/5, and mean clinical consistency 4.45/5. All twelve target-cluster outputs and all five patient-disjoint selected outputs passed. This establishes moderate-scale feasibility of the source-grounded, output-selected workflow in one target region; it does not establish cross-region generalization or make geometry a clinical-quality proxy. The next required test is the same preregistered workflow in a clinically distinct sparse region. Cluster 11 is frozen as that second region with all 19 available Tier-1 candidates, including eight patient-disjoint rows, under the same five-field transition-note contract.

The cluster-11 replication preserved final-output geometry (11/17 frozen anchors yielded a cluster-11 selected output) but did not reproduce the cluster-25 clinical result: blinded review passed 12/15 selected notes, with two unsupported medication additions and one critical omission of ledger-supported medications. Required section headings therefore improve structural completeness but cannot establish medication reconciliation. The combined cross-region result supports output-space targeting as a geometry mechanism, while showing that raw fact-only generation is not yet a uniformly clinically reliable enrichment method. Further raw scaling is stopped. The next phase is a medication-reconciliation detector evaluated against the completed blinded labels; only a detector with demonstrated sensitivity and acceptable rejection burden may become a prospective gate. A local LLM-as-judge is a subsequent option for unresolved semantic cases, whereas an editor is evaluated only after detection is reliable and must be compared against the unedited output.

The first medication-reconciliation diagnostic tested a conservative lexical comparison between reviewed ledger medication terms and generated medication sections across the completed cluster-25 and cluster-11 reviews. It caught only one of two human-labeled medication-addition failures while flagging many clinically passing notes, so it is not used as a rejection gate. This establishes that the remaining error is semantic reconciliation, not merely missing headings or exact medication-string coverage. Phase 3a therefore begins as an \\textbf{LLM-judge feasibility and calibration study}, not as a validated automatic gate. An approved independent local instruction-following model receives the compact verified medication ledger and generated note and must return evidence-first structured JSON: fact-to-note alignments, missing ledger medications, unsupported additions, contradictions, uncertainty, and a deterministic rejection recommendation. Cluster 25 is development-only for prompt, model, and threshold selection: in addition to the 20 completed selected-output reviews, a blinded 40-candidate development pack is sampled from the existing moderate-scale candidate pool after excluding those 20 outputs. Cluster 11 remains locked for one frozen leave-region-out evaluation. The current 20/15 split has only three severe cluster-11 failures, so even perfect detection is encouraging but statistically imprecise. Before any automatic prospective use, a fresh third region must test frozen judge settings with human review of selected outputs, judge-rejected outputs, and a random sample of judge-passing nonselected candidates. An editor remains downstream and may only rewrite judge-flagged notes under the same ledger, with an explicit comparison against the unedited output.

The completed 40-candidate cluster-25 calibration review found two severe medication failures among 40 additional candidates: a ledger-supported lactulose omission and an action contradiction that continued a ledger-stopped treatment. Both arose in candidates already rejected by the deterministic eligibility screen; this is compatible with, but does not prove, risk reduction because the sample is stratified and anchor-correlated. The 40 cases are combined with the 20 previously reviewed cluster-25 outputs for development only. The immediate task is model provision and calibration, not more note generation: select an approved independent local instruction-following judge using only cluster-25 development labels, freeze prompt, schema, and thresholds, and run the locked cluster-11 evaluation once with repeated deterministic inference. A new independent third region is required before the judge can influence prospective candidate selection.

\noindent\textbf{Revised mechanism and claim boundary.}
The project no longer claims that an inference-time edited embedding is reliably decoded into a faithful sparse-region note. Instead, embedding-space analysis identifies under-covered real-note regions and selects held-out anchors; a verified compact fact ledger supports local fact-conditioned synthesis; blinded factual review gates note validity; and only clinically passing notes are re-embedded to audit whether the final output enriches the intended real manifold. This preserves the cohort-enrichment objective while separating anchor selection, factual control, and output-space evaluation.

\begin{table}[h]
\centering
\small
\begin{tabularx}{\textwidth}{l l X}
\toprule
\textbf{Stage} & \textbf{Status} & \textbf{Next action} \\
\midrule
MIMIC-IV preprocessing + HADM alignment & \textbf{Done} & --- \\
Sentence embeddings ($\sim$332k) & \textbf{Done} & --- \\
HF datasets + long-sequence filter & \textbf{Done} & --- \\
ELM training (\texttt{checkpoint-8215}) & \textbf{Done} & --- \\
Leakage audit & \textbf{Done} & Report in all evaluations \\
Filtered-aligned leakage manifest refresh & \textbf{Done} & Use canonical \texttt{split\_manifest\_note\_level.csv} for all official joins \\
Vanilla held-out generation & \textbf{Done} & Official baseline established on 32,843 held-out rows \\
Structured manifest + vanilla audit & \textbf{Done} & Baseline audit PASS; use as official Phase 1 anchor \\
Whole-filtered-real coverage mapping & \textbf{Done} & Use all filtered train/dev/test embeddings to identify sparse/dense regions and candidate bias structure \\
Held-out real-only coverage baseline & \textbf{Done} & Preserve held-out baseline for later real-vs-synthetic validation \\
Real-vs-synthetic coverage mapping & \textbf{Done for vanilla baseline} & Use the held-out vanilla result as the comparison anchor for later pilots \\
Subgroup metadata build & \textbf{Done} & Join subgroup fields into future coverage and CAV reports \\
First global-axis CAV pilots + controls & \textbf{Done} & Keep as diagnostic evidence, not yet a scale-up mechanism \\
Inference-time steering + decoder adaptation & \textbf{Frozen negative/partial} & Report decoder compatibility as the limiting mechanism; do not add another heuristic steering family \\
Closed-loop output-space enrichment & \textbf{Positive basin-landing feasibility; clinical gate failed} & Retain the 106-note result and same-anchor 8192-token control as the geometry baseline; do not claim clinically usable enrichment yet \\
Source-grounded decoder rescue & \textbf{Prospective 30-case replication preparation} & The 4/4 fact-only smoke selects the method but has wide uncertainty; retain raw ELM baseline, freeze draft-correction arms, and test a prespecified replication \\
LLM editor / external reference / NER & Phase 3 & Start only after clinically credible closed-loop enrichment or a validated quality-control/editor strategy \\
\bottomrule
\end{tabularx}
\caption{Compact project status and next-action tracker.}
\end{table}

%=============================================================================
\section{Stage success criteria}
%=============================================================================

\textbf{Stage 1 success.}
\begin{itemize}[nosep]
  \item generated notes are readable, non-empty, and mostly non-collapsed,
  \item faithfulness is stable on both the full held-out set and the patient-disjoint subset,
  \item no major memorization or PHI red flags appear in baseline audit.
\end{itemize}

\textbf{Stage 2 success.}
\begin{itemize}[nosep]
  \item the full real manifold is mapped,
  \item under-covered clusters and subgroup skews are identifiable,
  \item vanilla coverage gaps are measurable.
\end{itemize}

\textbf{Stage 3 success.}
\begin{itemize}[nosep]
  \item output-space selection improves target-basin enrichment beyond a same-anchor, same-decoding vanilla control,
  \item accepted notes preserve source alignment, quality, diversity, and privacy-risk screens,
  \item blinded clinical and source-faithfulness review does not reveal unacceptable semantic degradation,
  \item acceptance rate and all rejected candidates are reported transparently.
\end{itemize}

\textbf{Stage 3 near-term redesign checklist.}
\begin{itemize}[nosep]
  \item Freeze input-space steering, projection, and decoder-adaptation families as negative/partial results.
  \item Retain the completed one-draw, default-prompt 8192-token matched vanilla control and paired review as the fair reference.
  \item Use a reviewed compact fact ledger without source spans as the sole factual authority for a local two-arm decoder smoke test: raw-draft correction and fact-only generation.
  \item Compare the untouched local backbone and \texttt{checkpoint-8215}; record model condition, ledger hash, raw-draft provenance, and token stopping metadata. The four-case smoke completed with 16 non-empty EOS-terminated B/C outputs and a blinded five-condition review pack.
  \item The blinded result selected \texttt{checkpoint-8215} fact-only generation: 4/4 passes, no unsupported major claims, and no critical omissions. This is a method-selection signal with wide four-case uncertainty, not evidence of near-perfect performance. Raw ELM and both draft-correction arms were 0/4 passes, so the untrusted draft is not retained as conditioning context.
  \item Run a prospective 30-case replication of the selected fact-only arm with raw ELM baseline, reviewer-approved compact ledgers, and blinded factual review. A nested 15-case untouched-backbone control is frozen before generation. Gate any further expansion on replication of factual preservation, unsupported-claim rate, critical omission rate, and clinically usable section coverage; evaluate geometry only after that gate passes.
\end{itemize}

\textbf{Replication decision rule.}
The 30 primary outputs and 15 nested controls completed blinded review. The primary \texttt{checkpoint-8215} fact-only arm met the predeclared factual gate (30/30 passes; zero unsupported-major-claim and critical-omission failures; all seven patient-disjoint cases passed), so post-review BGE re-embedding is approved as a pilot. The matched raw ELM baseline was 0/30 passes. The untouched-backbone subset was 15/15 passes, so the current result supports verified-fact-conditioned generation but does not establish a checkpoint advantage over the untouched backbone. Geometry analysis remains a 30-case pilot and must not be presented as cohort-wide enrichment without replication.

\textbf{Stage 4 success.}
\begin{itemize}[nosep]
  \item edited notes improve quality without breaking source alignment,
  \item synthetic-note augmentation improves at least one downstream clinical NLP task.
\end{itemize}

%=============================================================================
\section{Appendix: implementation notes}
%=============================================================================

\small
This appendix preserves technical detail for the engineering team without front-loading the clinical narrative.

\subsection{Key implementation notes}

\begin{itemize}[nosep]
  \item Each manifest row should preserve source ids, split, generation condition, decoding settings, leakage flags, and generated text.
  \item The source manifest should preserve stable embedding/source provenance for the full cohort.
  \item Recommended \texttt{generation\_condition} values: \texttt{vanilla}, \texttt{multiseed\_vanilla}, \texttt{random\_shift}, \texttt{cav\_steered}, \texttt{elm\_edited}.
\end{itemize}

\subsection{ELM implementation}

\begin{itemize}[nosep]
  \item Encoder: \texttt{BAAI/bge-large-en-v1.5} (1024-d)
  \item Decoder: Llama-3.1-8B with adapter injection (\texttt{open-elm/src/model.py})
  \item Standard generation input: prepared HF dataset \texttt{encoded\_testing\_filtered}, not raw \texttt{.npy} alone
  \item Practical implication for Phase 2b: the steering operator should respect the adapter's training distribution; unconstrained off-manifold edits are less likely to decode into the intended region
\end{itemize}

\subsection{Manuscript figures (target)}

Pipeline diagram; leakage/overlap table; vanilla faithfulness table; whole-filtered-real manifold map; held-out real-vs.\ vanilla coverage map; CAV tradeoff curve; privacy audit table; exemplar notes (real vs.\ vanilla vs.\ steered).

%=============================================================================
\section{Conclusion}
%=============================================================================

This project is a staged framework for source-conditioned synthetic discharge-summary generation. Phase 1 builds the audited vanilla baseline and the real-manifold map. Phase 2 tests whether targeted steering can fill under-covered regions better than matched controls. Phase 3 improves note quality if needed and tests whether the resulting synthetic corpus helps real clinical NLP tasks.

\begin{thebibliography}{99}
\bibitem{xu2021synthetic}
Xu, J., Liu, T., Pakhomov, S., et al.
\newblock Are synthetic clinical notes useful for real natural language processing tasks: a case study on clinical entity recognition.
\newblock \emph{JAMIA}, 28(10):2193--2201, 2021.

\bibitem{melamud2019towards}
Melamud, O. and Shivade, C.
\newblock Towards automatic generation of shareable synthetic clinical notes using neural language models.
\newblock \emph{Clinical NLP Workshop}, 2019.

\bibitem{kim2018interpretability}
Kim, B., Wattenberg, M., Gilmer, J., et al.
\newblock Interpretability beyond feature attribution: quantitative testing with concept activation vectors (TCAV).
\newblock \emph{ICML}, 2018.

\bibitem{morris2023vec2text}
Morris, J. X., Kuleshov, V., Shmatikov, V., and Rush, A.
\newblock Text Embeddings Reveal (Almost) As Much As Text.
\newblock \emph{arXiv:2310.06816}, 2023.
\newblock \url{https://arxiv.org/abs/2310.06816}

\bibitem{bartoszcze2025repe}
Bartoszcze, L., Munshi, S., Sukidi, B., et al.
\newblock Representation Engineering for Large-Language Models: Survey and Research Challenges.
\newblock \emph{arXiv:2502.17601}, 2025.
\newblock \url{https://arxiv.org/abs/2502.17601}

\bibitem{nguyen2025matsteer}
Nguyen, D., Prasad, A., Stengel-Eskin, E., and Bansal, M.
\newblock Multi-Attribute Steering of Language Models via Targeted Intervention.
\newblock \emph{ACL}, 2025.
\newblock \url{https://aclanthology.org/2025.acl-long.1007/}

\bibitem{vanderweij2025cast}
van der Weij, T., et al.
\newblock Mitigating Content Effects on Reasoning in Language Models Through Fine-Grained Activation Steering.
\newblock \emph{arXiv:2505.12189}, 2025.
\newblock \url{https://arxiv.org/abs/2505.12189}

\bibitem{carlini2021extracting}
Carlini, N., Tramer, F., Wallace, E., et al.
\newblock Extracting training data from large language models.
\newblock \emph{USENIX Security}, 2021.
\end{thebibliography}

\end{document}
